% preamble-exam-zh.tex — 共享 preamble
% 用于所有 exam-zh 模板

\documentclass{exam-zh}

% ---------- 基础配置 ----------
\examsetup{
  page/size = a4paper,
  page/show-head = false,
  page/show-foot = false,
  sealline/show = false,
  font = times,
  math-font = xits,
}

\usepackage{tcolorbox}
\usepackage{needspace}
\tcbuselibrary{skins,breakable}

% ---------- 自定义教学环境 ----------

% 原题卡片 (edu-problem-card)
\newtcolorbox{problemcard}{
  colback=black!2,
  colframe=black!80,
  boxrule=1.5pt,
  arc=0pt,
  left=3mm, right=3mm, top=3mm, bottom=3mm,
}

% 解题路线图 (edu-route)
\newtcolorbox{routemap}{
  colback=black!2,
  colframe=black!60,
  boxrule=1pt,
  arc=0pt,
  left=3mm, right=3mm, top=3mm, bottom=3mm,
}

% 关键想法 (edu-key-idea)
\newtcolorbox{keyideabox}{
  colback=yellow!3,
  colframe=black!80,
  boxrule=1.5pt,
  arc=0pt,
  left=3mm, right=3mm, top=3mm, bottom=3mm,
}

% 易错提醒 (edu-mistake)
\newtcolorbox{mistakebox}{
  colback=red!3,
  colframe=black!30,
  boxrule=1pt,
  arc=0pt,
  borderline west={3pt}{0pt}{red!80!black},
  left=3mm, right=3mm, top=2.5mm, bottom=2.5mm,
}

% 提示 (edu-hint)
\newtcolorbox{hintbox}{
  colback=black!2,
  colframe=black!30,
  boxrule=1pt,
  arc=0pt,
  dashed,
  left=3mm, right=3mm, top=2mm, bottom=2mm,
}

% 思考题 (edu-question)
\newtcolorbox{thinkbox}{
  colback=white,
  colframe=black!60,
  boxrule=1pt,
  arc=0pt,
  dashed,
  left=2.5mm, right=2.5mm, top=2.5mm, bottom=2.5mm,
}

% 教师备注 (edu-teacher-note) — teacher 版渲染
\newtcolorbox{teachernote}{
  colback=black!3,
  colframe=black!30,
  boxrule=1pt,
  arc=0pt,
  left=2.5mm, right=2.5mm, top=2.5mm, bottom=2.5mm,
  fontupper=\small,
  colupper=black!50,
}

% 训练目标 (edu-training-goal) — teacher 版渲染
\newtcolorbox{traininggoal}{
  colback=white,
  colframe=black!30,
  boxrule=0.5pt,
  arc=0pt,
  left=2mm, right=2mm, top=1mm, bottom=1mm,
  fontupper=\small,
}

% 答题区 (edu-answer-lines)
\newcommand{\answerarea}[1][40mm]{%
  \vspace{3mm}%
  \begin{tcolorbox}[
    colback=white,
    colframe=black!30,
    boxrule=1pt,
    arc=0pt,
    height=#1,
    valign=top,
  ]
  \end{tcolorbox}%
}

% ---------- 字体设置 ----------
\setCJKmainfont{SimSun}[AutoFakeBold, AutoFakeSlant]
\setCJKsansfont{SimHei}
\setCJKmonofont{FangSong}

% ---------- 页面设置 ----------
\usepackage[top=16mm, bottom=16mm, left=14mm, right=14mm]{geometry}

\examsetup{
  question/show-answer = false,
  fillin/show-answer = false,
  solution/show-solution = hide,
}

\begin{document}

\title{一次函数专题练习}
\subject{%
  图像与解析式 · 八年级 · 数学}
\maketitle

\section*{一、核心练习}


\begin{question}[points=4]
  函数 $y=2x-1$ 的图像经过下列哪个点？
  \begin{choices}
    \item $(0,1)$
    \item $(1,1)$
    \item $(1,0)$
    \item $(-1,1)$
  \end{choices}
\end{question}




\begin{question}[points=4]
  一次函数 $y=-x+3$ 与 $x$ 轴的交点坐标是 \underline{\hspace{2cm}}。
  \fillin
\end{question}




\begin{problem}[points=8]
  已知一次函数 $y=kx+b$ 的图像经过点 $(2,5)$ 和点 $(-1,-1)$，求 $k$ 和 $b$ 的值。
  \answerarea[45mm]
\end{problem}




\begin{question}[points=4]
  一次函数 $y=-2x+4$ 的图像不经过第 \underline{\hspace{1cm}} 象限。
  \begin{choices}
    \item 一
    \item 二
    \item 三
    \item 四
  \end{choices}
\end{question}



\clearpage
\section*{二、参考答案}


\begin{keyideabox}
  \textbf{第 1 题}
  选 B。代入 $x=1$，得 $y=2\times1-1=1$，即点 $(1,1)$ 在图像上。
\end{keyideabox}




\begin{keyideabox}
  \textbf{第 2 题}
  $(3,0)$。令 $y=0$，则 $-x+3=0$，$x=3$。
\end{keyideabox}




\begin{keyideabox}
  \textbf{第 3 题}
  $k=2, b=1$。由方程组 $\begin{cases}2k+b=5\\-k+b=-1\end{cases}$ 解得。
\end{keyideabox}




\begin{keyideabox}
  \textbf{第 4 题}
  选 C。$k=-2<0, b=4>0$，图像过一、二、四象限。
\end{keyideabox}




\end{document}
